\section{Multiple points and exhaustive search for pseudolines}\label{sec:enum}

We now turn from straight line arrangements to the general case of pseudoline arrangements.
An \emph{arrangement of $n$ pseudolines} is a collection of $n$ unbounded simple curves in the plane in which every two meet exactly once and cross there.
Every line arrangement is a pseudoline arrangement, so the maximum number $K(n)$ of triangular faces in an arrangement of $n$ lines does not exceed its pseudoline analogue $K_{\mathrm{ps}}(n)$:
\begin{equation}\label{eq:K-ps}
K(n)\le K_{\mathrm{ps}}(n)\overset{\mathrm{def}}{=}\max\{a_3(\mathcal A):\ \mathcal A\ \text{is an arrangement of } n \text{ pseudolines}\}.
\end{equation}

The only known upper bound for arrangements with multiple points is \eqref{eq:a3-bound}, and among the known even arrangements it is attained only for $n=4$.
We wrote a program that enumerates all arrangements with at least a given number $K$ of triangles, with vertices of any multiplicity.
For $n=6,8,10,12$ the enumeration finishes in reasonable time, and for $n=14$ the program shows that no arrangement has 55 or more triangles.
Together with line arrangements attaining them, this gives the exact values $K(n)=K_{\mathrm{ps}}(n)=7,15,25,38,54$, which equal the conjectured bound \eqref{eq:conj-bound}.
The program does not use the parity of $n$, and for $n=11$ it rules out 33 triangles, so $K(11)=32$.

\subsection{Representing arrangements by words}

The search works with a combinatorial representation of arrangements obtained by \emph{sweeping}.
Take a curve that meets every pseudoline exactly once and move it through the arrangement so that it passes the intersection points, or \emph{crossings}, one at a time.
The pseudolines cross the curve in some order, which is a permutation.
Passing a crossing changes this permutation, and the sequence of these changes encodes the arrangement.

A straight line sweeps any line arrangement, but for pseudolines a sweep has to be shown to exist.
By Goodman's classical theorem, every pseudoline arrangement is isomorphic to a wiring diagram, that is, admits such a sweep~\cite{goodman80}, \cite[Theorem~5.1.4]{fg17}.
However, the theorem concerns the projective plane: the isomorphism need not preserve the line at infinity, and with it the bounded faces and the number $a_3$.
We therefore use the sweeping theorem of Snoeyink and Hershberger~\cite[Theorem~3.1]{sh89}: every arrangement in the affine plane, simple or not, can be swept by a pseudoline.
Such a sweep takes place in the affine plane, so the wiring diagram it produces preserves the partition of the faces into bounded and unbounded ones.

We have already considered simple arrangements in detail in~\cite[Subsection~2.1]{pu26b}, where a sweep gives a reduced word for the longest permutation $w_0$, and each letter exchanges two adjacent pseudolines.
At a crossing of multiplicity $m$, the permutations of the allowable sequence~\cite{gp93} before and after it differ by the reversal of a block of $m$ consecutive elements instead of two.
We call the positions in these permutations \emph{slots} and number them from $0$ to $n-1$.
A pseudoline arrangement thus corresponds to a word in the letters
\begin{equation}\label{eq:letter}
(g,m):\qquad \text{reversal of the block in slots } g,\ldots,g+m-1,
\end{equation}
with one letter for each crossing.
A letter $(g,m)$ accounts for the $\binom m2$ pairs of pseudolines that cross at one point.

A \emph{word} is any sequence of letters, and a \emph{wiring diagram} is a drawing of a word. We call its curves \emph{wires}.
The letters act on the current permutation from left to right, starting with the identity permutation.
A word is \emph{partial} if no two wires cross twice in it, and \emph{complete} if every two wires cross exactly once.
Appending letters $V$ to a partial word $W$ gives an \emph{extension} $WV$ of $W$, and $V$ is a \emph{continuation} of $W$.
A complete extension is a \emph{completion} of $W$.

The verdicts of the search need only a word for every arrangement, and the Snoeyink--Hershberger sweep~\cite{sh89} provides one: if the search emits no words, no arrangement of the kind sought exists.
In the other direction, every emitted word gives a wiring diagram, which is an example of such an arrangement.

Two consecutive letters $(g,m)$ and $(g',m')$ of a partial word commute when they act on disjoint sets of wires, that is, when their blocks of slots do not overlap:
\begin{equation}\label{eq:mult-commut}
(g,m)(g',m') = (g',m')(g,m)\iff g'\geq g+m \lor g \geq g'+m'.
\end{equation}
A sweep is free to take commuting letters in either order, so we regard the drawings of words that differ by such swaps as the same wiring diagram.
The search removes this commutation freedom and emits one word from each commutation class.

\begin{lemma}[canonical form]\label{lem:canonical}
Each commutation class contains exactly one word with no pair of consecutive letters $(g,m)(g',m')$ such that $g'+m'\le g$, that is, with the block of the second letter entirely to the left of the block of the first. This word is the \emph{canonical form} of the class.
\end{lemma}

\begin{proof}
The proof follows~\cite[Lemma~4.7]{pu26b}. Swapping consecutive commuting letters \eqref{eq:mult-commut} either removes a forbidden pair or creates one. Remove forbidden pairs one at a time. Each such swap decreases by one the number of pairs of letters in which the later block lies entirely to the left of the earlier one, so the process terminates. Two forbidden pairs overlap only in a triple $abc$ where the block of $c$ is to the left of the block of $b$, which in turn is to the left of the block of $a$. Then the block of $c$ is to the left of the block of $a$, and either order of swaps leads to $cba$. By Newman's lemma~\cite{newman42}, all words of the class reach the same result.
\end{proof}

The search emits only canonical forms: after a letter $(g,m)$ it forbids the letters $(g',m')$ with $g'\le g-m'$.

One more freedom comes from the sweep rather than from the arrangement: the choice of the start, or ``direction''.
The start can be chosen so that the first letter $(g_0,m_0)$ begins in slot $g_0=0$.
This choice, the \emph{first-letter fixation}, is useful for negative verdicts, where it suffices to search over words with first letter $(0,m)$.

\begin{lemma}[first letter $g_0=0$]\label{lem:first-letter}
Every affine arrangement of $n\ge2$ pseudolines, with or without multiple points, has a word whose first letter is $(0,m)$ for some $2\le m\le n$.
The canonical form (\cref{lem:canonical}) of this word also begins with $(0,m)$.
\end{lemma}

\begin{proof}
Take any sweep of the arrangement $\mathcal A$.
Let $(g,m)$ be its first letter and $v$ the crossing it produces.
The sweep passes no crossings before $v$, so the initial curve can be pushed to infinity without touching any crossing. There it runs along $n$ consecutive ends of pseudolines, and the order of the wires on it is the order of these ends.

Any $n$ consecutive ends contain one end of each pseudoline, since the two ends of a pseudoline are separated by $n-1$ ends on each side.
Rotate the initial curve past $g$ ends.
The new curve $\gamma$ still crosses each pseudoline once, and its initial order is a cyclic shift of the old one, with the wires of $v$ in slots $0,\ldots,m-1$.
The two positions of the initial curve differ only near the ends of the $g$ transferred pseudolines, and none of these pseudolines passes through $v$.
Hence the wires of $v$ still have no other crossings between $\gamma$ and $v$, and $\gamma$ can be moved across $v$ to a curve $\gamma'$. This move is the first letter $(0,m)$.

Since $\gamma'$ crosses each pseudoline exactly once, $\mathcal A\cup\{\gamma'\}$ is a pseudoline arrangement, and by~\cite[Theorem~3.1]{sh89} it can be swept starting from $\gamma'$ towards the side of $\gamma'$ that contains all crossings of $\mathcal A$ other than $v$.
Such a sweep uses two kinds of operations: passing an empty triangle, and passing or taking the first ray~\cite[Lemma~3.1]{sh89}.
The curve $\gamma'$ already crosses all pseudolines, so it never takes a ray, and the crossings of $\mathcal A$ are passed by the first operation.
A triangle in~\cite{sh89} is not a face of the arrangement but the region bounded by the sweep curve and two wires consecutive on it, up to their common crossing.
A crossing of multiplicity $m$ has $m-1$ such regions, and the sweep passes all of them at once, with one letter~\cite[p.~356]{sh89}.
The result is a complete word with first letter $(0,m)$.

Bringing this word to canonical form removes forbidden pairs as in the proof of \cref{lem:canonical}.
The first letter stays in place: moving it would require the block of the second letter to lie entirely to the left of slot 0, which is impossible.
\end{proof}

\subsection{Counting faces}

Euler's formula gives budget constraints for the search.

\begin{lemma}\label{lem:faces}
Let $\mathcal A$ be an arrangement of $n\ge3$ pseudolines.
Then the number of its bounded faces is
\begin{equation}\label{eq:euler-F}
F_b(\mathcal A)=\binom{n-1}2-\sum_v\binom{m(v)-1}2,
\end{equation}
where the sum runs over all vertices $v$ and $m(v)$ is the multiplicity of $v$.
\end{lemma}

\begin{proof}
The $k_\ell$ vertices on a pseudoline $\ell$ divide it into $k_\ell-1$ bounded segments and two rays.
Each vertex $v$ lies on $m(v)$ pseudolines, so counting the incident pairs ``vertex, pseudoline'' in two ways gives $\sum_\ell k_\ell=\sum_v m(v)$. The number of bounded edges is therefore
\begin{equation}\label{eq:edges}
E_b=\sum_\ell(k_\ell-1)=\sum_v m(v)-n.
\end{equation}
The graph of vertices and bounded edges is connected: vertices $u\in\ell$ and $v\in\ell'$ are joined along $\ell$ and $\ell'$ through their common point.
The rays lie in the unbounded face of this graph, so its bounded faces are the bounded faces of $\mathcal A$, and by Euler's formula for a connected plane graph their number is $F_b=E_b-V+1$, where $V$ is the number of vertices.
Since every pair of pseudolines crosses exactly once, the total number of pairwise crossings is
\begin{equation}\label{eq:pair-cross}
P(\mathcal{A})=\sum_v\binom{m(v)}2=\binom n2,
\end{equation}
and $\binom m2-\binom{m-1}2=m-1$ gives
\[
\sum_v\bigl(m(v)-1\bigr)=\binom n2-\sum_v\binom{m(v)-1}2 .
\]
Substituting this into $E_b-V+1=\sum_v(m(v)-1)-n+1$ and using $\binom n2-n+1=\binom{n-1}2$, we obtain \eqref{eq:euler-F}.
\end{proof}

For a simple arrangement every vertex has $m(v)=2$, all terms of the sum vanish, and \eqref{eq:euler-F} gives $\binom{n-1}2$, the ceiling on the number of bounded faces of any arrangement. The sum shows how many faces multiple points take away: a vertex of multiplicity $m$ destroys $\binom{m-1}2$ bounded faces. Not all of the remaining bounded faces are triangles. Let $a_k(\mathcal A)$ be the number of bounded faces with $k$ sides. Bounded digons do not occur, since two pseudolines cross only once. Then
\begin{equation}\label{eq:a3-split}
a_3(\mathcal A)=\binom{n-1}2-\sum_v\binom{m(v)-1}2-(a_4+a_5+\ldots).
\end{equation}
Every multiple point and every bounded non-triangular face pushes the number of triangles below the ceiling $\binom{n-1}2$, and the target $a_3\ge K$ tolerates a shortfall of at most $\binom{n-1}2-K$ triangles.

\begin{definition}\label{def:budget}
We call $B=\binom{n-1}2-K$ the \emph{budget} of the target $K$ and the number
\begin{equation}\label{eq:spent}
S(\mathcal A)=\sum_v\binom{m(v)-1}2+a_4+a_5+\ldots
\end{equation}
the \emph{spending} of the arrangement: a vertex of multiplicity $m$ spends $\binom{m-1}2$ units of the budget, and each bounded non-triangular face spends one.
\end{definition}

\begin{corollary}\label{cor:budget}
Every arrangement satisfies $a_3(\mathcal A)=\binom{n-1}2-S(\mathcal A)$, so $a_3(\mathcal A)\ge K$ if and only if
\begin{equation}\label{eq:budget}
S(\mathcal A)\le B .
\end{equation}
\end{corollary}

The budget is shared between multiple points and non-triangular faces.
For instance, for $n=8$ and $K=16$ we have $B=5$; a vertex of multiplicity 4 takes 3 units, leaving 2 for all non-triangular faces of the arrangement.

A triple point costs only one unit of the budget, so it can pay for itself: for the price of one face and three segments, some of which may be unused, it lets neighboring triangles share a segment.\footnote{The condition $a_3\ge K$ can be written as two budget constraints, one from counting faces \eqref{eq:budget} and one from counting segments, $w\le b$ in the notation of \eqref{eq:waste-limit} and \eqref{eq:segm-balance}.
For an arrangement, that is, for a complete word, they are equivalent, but on a partial word they give different prunings, so the search uses both (see \cref{sec:segment-count}).}
Empirically, the cost of lost faces grows with the multiplicity faster than the possible gain, and the optima contain no points of higher multiplicity (see \cref{tab:triple-points}).

\subsection{Search algorithm}\label{sec:search}

The search builds a word from left to right, one letter per step, and
emits exactly the canonical forms with $a_3\ge K$. Define it as the following recursive
procedure:

\begin{verbatim}
Function Search(W), where W is a partial word:
   if S(W) > B, return;                    // a3 >= K is unreachable
   if P(W) = C(n,2), output W and return;  // W is complete
   else for each letter (g,m) admissible at W, call Search(W·(g,m)).
\end{verbatim}

\noindent
The search starts from the empty word.
Let $a[0],\ldots,a[n-1]$ be the permutation\footnote{Square brackets index slots.
In~\cite[Section~4]{pu26b} this permutation was written $a_i$;
the notation is changed to avoid confusion with the triangle count $a_3$.}
of wire labels obtained from the identity permutation by applying $W$; that is, $a[i]$
is the wire occupying slot $i$ after the crossings of $W$.
A letter $(g,m)$ reverses the block $a[g],\ldots,a[g+m-1]$.
The number of wire pairs already crossed in $W$ is
$P(W)=\sum_v\binom{m(v)}2$, as in \eqref{eq:pair-cross}.
The word $W$ is complete when $P(W)=\binom n2$, and then $a[i]=n-1-i$ is the longest permutation $w_0$.
The spending $S(W)$ is defined by \eqref{eq:spent} over the vertices of $W$
and the bounded faces that $W$ has already closed.

A letter $(g,m)$ is \emph{admissible} at $W$ if it satisfies all of the following conditions.

(i) $0\le g\le n-m$ and $a[g]<a[g+1]<\cdots<a[g+m-1]$: no two wires of the block have crossed yet, so the letter crosses each pair of them exactly once and produces all $\binom m2$ pairwise crossings at once.

(ii) $S(W)+\binom{m-1}2\le B$: the new vertex of multiplicity $m$ produced by $(g,m)$ destroys $\binom{m-1}2$ bounded faces (\cref{lem:faces}), and its cost, added to the current spending, has to fit into the budget \eqref{eq:budget}.
The same condition automatically bounds $m$, so the exhaustive search needs no separate cap on the multiplicity.

(iii) If $W$ is nonempty and its last letter is $(\widehat g,\widehat m)$, then $g\ge \widehat g-m+1$. Otherwise the block of $(g,m)$ lies entirely to the left of the block of $(\widehat g,\widehat m)$, and $W\cdot(g,m)$ contains a pair forbidden by \cref{lem:canonical}, so none of its extensions is canonical.

Condition (i) is responsible for correctness. Without it, the search would emit words that correspond to no
wiring diagram, since some pair of wires would cross twice.
Condition (iii) leaves exactly one word, the canonical form, from each commutation class (\cref{lem:canonical}).
Condition (ii) is the \emph{budget pruning} (P1). The letter adds at least $\binom{m-1}2$ to the spending, and appending letters to a partial word $W$ can only increase $S(W)$. So once $S(W)>B$, every completion of $W$ contradicts \cref{cor:budget}, and the discarded branch contains no words with $a_3\ge K$.

To count the second term of the spending, the search has to know the number of sides of each face.
A bounded face of a wiring diagram receives its vertices from letters.
We call the current face between the wires in slots $i$ and $i+1$ \emph{face $i$}.
A letter that crosses the wires in slots $i$ and $i+1$, that is, whose block contains both slots, closes face $i$ and opens a new face $i$.
Letters whose block ends at slot $i$ or begins at slot $i+1$ touch only one of these wires and give face $i$ intermediate vertices.
So the number of vertices of every face, and in particular the numbers of triangles and non-triangular faces, can be counted directly from the partial word $W$.

\subsection{Prunings from properties of wiring diagrams}

Conditions (i)--(iii) alone do not yet give a practical exhaustive search.

\begin{lemma}[face boundedness]\label{lem:bounded-face}
If $\min(a[0],\ldots,a[i])<n-1-i$, that is, the set of wire labels in slots $0,\ldots,i$ is not yet final, then face $i$ closes in every completion of the partial word $W$. Otherwise it closes in no completion.
In particular, a face $i$ already opened by some letter is bounded if and only if $\min(a[0],\ldots,a[i])<n-1-i$.
\end{lemma}

\begin{proof}
A letter whose block does not contain both slots $i,i+1$ permutes wires within slots $0,\ldots,i$ or within slots $i+1,\ldots,n-1$ and does not change the set of labels on the left. Only a letter that closes face $i$ can change this set.
In the final permutation, slots $0,\ldots,i$ hold the labels $n-1,\ldots,n-1-i$.
If the labels on the left are not yet final, their set changes in every completion, that is, face $i$ closes.
If they are final, then every label on the left exceeds every label on the right.
In particular, $a[i]>a[i+1]$, now and after every letter that does not close face $i$, so by condition (i) face $i$ never closes (cf.~\cite[Remark~B.8]{pu26b}).
\end{proof}

The set of labels $\{a[0],\ldots,a[i]\}$ changes only when face $i$ closes, so whether this face is bounded is decided at the moment it opens.
A bounded face $i$ that already has at least two intermediate vertices will have at least four vertices when it closes.
Let $\delta S(W)$ be the number of such faces.
They are not yet counted in $S(W)$, and each of them adds one to $S$ when it closes.
This gives a pruning by \emph{early non-triangle count} (P2): the search discards a branch as soon as $S(W)+\delta S(W)>B$, since by \cref{cor:budget} the branch contains no words with $a_3\ge K$.

\emph{Forced descent of a suffix-maximal wire} (P3). In~\cite[Appendix~B.1]{pu26b} we described a pruning based on the forced descent of a wire whose label exceeds the labels of all wires to its right in the permutation. The idea carries over to multiple points, but now only the last slot of the block of the next letter is forced ($g+m-1=\widehat g$).

\begin{lemma}[forced descent]\label{lem:forced-descent}
Let $(\widehat g,\widehat m)$ be the last letter of a partial word $W$ in canonical form. Suppose the wire $w=a[\widehat g]$ is not yet in its final slot $n-1-w$, and all labels to its right in the permutation ($a[\widehat g+1],a[\widehat g+2],\ldots$) are smaller than $w$. Then in every canonical completion of $W$ the next letter has the form $(\widehat g-m+1,m)$, that is, its block ends at slot $\widehat g$.
\end{lemma}

\begin{proof}
The wire $w$ is not yet in its final slot, so some later letter moves it. Let $\alpha$ be the first such letter.
The block of $\alpha$ contains slot $\widehat g$ but not $\widehat g+1$: all labels to the right of $w$ are smaller than $w$, while by condition (i) the block of $\alpha$ must be increasing. Hence this block ends at $\widehat g$. The letters between $(\widehat g,\widehat m)$ and $\alpha$ do not touch slot $\widehat g$, so each of them lies entirely to the left or entirely to the right of it. A letter to the left of $\widehat g$ cannot come right after $(\widehat g,\widehat m)$, since the two would form a pair forbidden by \cref{lem:canonical}. So if there are letters in between, the first of them lies to the right of $\widehat g$, and the last letter on the right is followed by a letter to the left of $\widehat g$ or by $\alpha$ itself, and these two letters again form a forbidden pair. Hence $\alpha$ comes right after $(\widehat g,\widehat m)$.
\end{proof}

After $\alpha$ the wire $w$ occupies the initial slot of its block, and all labels to its right are again smaller than $w$. By induction, the lemma applies to this wire until it reaches its final slot. The multiplicity of the next letter remains arbitrary, so the letters of the descent chain are not determined in advance: for each possible multiplicity $m$ the search keeps a single block start $\widehat g-m+1$.

\subsection{Counting segments and the perturbation principle}\label{sec:segment-count}

In the introduction we wrote down the segment balance \eqref{eq:simple-segm} for the simple case. We treated this balance in more detail in~\cite[Subsection~2.3]{pu26b}.
Each of the $n(n-2)$ bounded segments is a side of at most one triangle,
since otherwise either simplicity or the requirement that two pseudolines cross only once would be violated (see~\cref{lem:double-segment}).
We now consider what changes in the presence of multiple points.

Roudneff~\cite{roudneff96} described a perturbation principle for multiple points.
Multiple points fall into two groups. Points of the first group admit a deformation of the pseudolines near the point that lowers its multiplicity and does not decrease the total number of triangles. For points of the second group, every such deformation decreases this number.
Only the latter are of interest, since only they can give more triangles than simple crossings.
All such points are reduced in the sense of the following definition.

\begin{definition}\label{def:reduced}
The \emph{sectors} of a vertex $v$ of multiplicity $m$ are the $2m$ faces incident to it.
A vertex of multiplicity $3$ or $4$ is \emph{reduced} if its sectors that are not bounded triangles form one of the types of~\cite[Fig.~3.5]{roudneff96}.
For $m=3$, each of the two triples of alternate sectors contains at most one such sector.
For $m=4$, there are at most two such sectors, and if there are two, they are opposite.
An arrangement is \emph{reduced} if all its vertices of multiplicity $\ge3$ are reduced.
\end{definition}

For example, under a deformation a triple point turns into a new triangle, and the faces in three alternate sectors each gain an extra side.
The direction of the deformation determines which triple of sectors gains them.
At a reduced triple point each triple contains at least two triangular sectors, so every direction of deformation destroys at least these two triangles and creates only one, and the total number of triangles decreases.

\begin{lemma}[perturbation]\label{lem:perturbation}
If a vertex $v$ of multiplicity $m\ge3$ of an arrangement $\mathcal A$ is not reduced, then there is an arrangement $\mathcal A'$ of the same number of pseudolines with $a_3(\mathcal A')\ge a_3(\mathcal A)$ and more vertices.
\end{lemma}

\begin{proof}
We carry over Roudneff's proof in the projective plane~\cite[\S3]{roudneff96} to the affine setting.
In a small disk $D$ around $v$, replace the $m$ arcs through $v$ by an arrangement $U$ of $m$ arcs with the same endpoints, any two of which cross once.
The faces that meet $\partial D$ correspond to the sectors of $v$ and do not change outside $D$. So boundedness is preserved, and the number of sides can only grow: a sector remains a triangle exactly when its local face in $U$ is a wedge of two rays with a common vertex.
The new faces lie in $D$ and are bounded.
Therefore $a_3$ does not decrease if $U$ has at least as many bounded triangles as there are destroyed triangular sectors, and the number of vertices grows, since $U$ has more than one vertex.

For $m=3,4,5$ take $U$ from~\cite[Fig.~3.3, 3.4, 3.1]{roudneff96}. For $m=5$ the number of triangles never decreases, and for $m=3,4$ it decreases only at vertices of the types of~\cite[Fig.~3.5]{roudneff96}, that is, at reduced vertices.
Roudneff's case analysis uses only which sectors are marked as non-triangular.
We apply it with the sectors that are not bounded triangles marked, as in \cref{def:reduced}; an unbounded sector is then never destroyed as a triangle, and the affine count is never worse.

For $m\ge6$ take the arrangement $\mathcal B_m$ of F\"uredi and Pal\'asti~\cite[Example~2, Property~3]{fp84}. Its number of projective triangles is $p_3(\mathcal B_m)=m(m-3)/3+6-2\varepsilon/3$, where $\varepsilon=0$ for $m\equiv0\pmod3$ and $\varepsilon=2$ otherwise, and $p_3(\mathcal B_m)\ge2m$ for all $m\ge6$ (with equality for $m=6,7$).
Cut the projective plane along a pseudoline in general position. Each triangle of $\mathcal B_m$ that this pseudoline crosses gives a wedge, and the triangles it misses remain bounded triangles.
Hence there are at least $2m$ wedges and triangles in total, and there are no more destroyed sectors than new triangles.
\end{proof}

Multiple points change the segment balance in two ways. At a vertex of multiplicity $m$, each of the $m$ pseudolines through it meets $m-1$ others at once and so loses $m-2$ bounded segments. Moreover, a bounded segment with a multiple endpoint can be a side of two triangles (\cref{lem:double-segment}). We call such segments \emph{double} and denote their number by $d(\mathcal{A})$. The segment balance (\cref{lem:segment-balance}) then reads
\begin{equation}\label{eq:non-simple-segm}
3a_3(\mathcal{A})+u(\mathcal{A})+3t_3(\mathcal{A})+8t_4(\mathcal{A})+\ldots-d(\mathcal{A})=n(n-2),
\end{equation}
where $t_m(\mathcal{A})$ is the number of vertices of multiplicity $m$, and $d(\mathcal{A})$ compensates for double segments counted twice in $3a_3(\mathcal{A})$.

This is how multiple points get around the bound \eqref{eq:blanc-even} for simple arrangements: the lost segments may be the unused ones that are unavoidable in the simple case, and double segments let triangles share sides. By \cref{lem:perturbation}, only reduced multiple points can gain triangles in these ways.

\emph{Segment pruning} (P4). A discharging argument~\cite{cw17} applied to the segment balance \eqref{eq:non-simple-segm} for crossings of arbitrary multiplicity leads to the inequality $u(\mathcal{A})\le n(n-2)-3K+\tfrac32t_3(\mathcal{A})$ (\cref{cor:unused-bound}). The search abandons a subtree as soon as this inequality can no longer hold.

\emph{State table pruning} (P5). A hash table stores the state of the search once a subtree has been traversed without finding any word with the target number of triangles $a_3\geq K$. The search skips a node whose state is \emph{dominated} by a stored one: it has the same permutation and the same $\widehat g$, and every completion of it gives at most as many triangles as the same completion of the stored word. Such nodes occur often: for $n=10$ the table alone reduces the search tree by a factor of about 400. In the search for perfect arrangements~\cite[Section~4]{pu26b}, a similar table removed few nodes and made each node more expensive, so we did not use it there.

\emph{Doomed-face count} (P6). \Cref{app:doomed} extends the early non-triangle count (P2) to all \emph{doomed} faces, that is, faces that can no longer become bounded triangles.

The justification of these three prunings is considerably longer and subtler than that of the others, and they are harder to implement.
We therefore give their details in \cref{app:prunings} and enable them in the search by a separate run parameter, so that numerical results independent of \cref{app:prunings} remain available.
The parameter does not change the output, as we checked on the exhaustive enumeration for $n=10$.
The theorems of this section for $n<12$ rest only on runs without these prunings.

\emph{Pruning of non-reduced arrangements} (P7). For a negative verdict, \cref{lem:perturbation} allows the search to consider only arrangements with the largest number of vertices.

\begin{corollary}\label{cor:reduced}
Among arrangements of $n$ pseudolines with $a_3\ge K$, every arrangement with the largest number of vertices is reduced.
It is represented by a word in canonical form with first letter $(0,2)$.
\end{corollary}

\begin{proof}
If such an arrangement were not reduced, \cref{lem:perturbation} would give an arrangement with $a_3\ge K$ and more vertices.
By \cref{lem:first-letter}, the arrangement has a word with first letter $(0,m)$, and the canonical form of this word also begins with $(0,m)$.
The vertex produced by the first letter is the first one on each of its $m$ pseudolines, so $m$ consecutive rays meet at it, and $m+1$ unbounded sectors are incident to them.
Hence this vertex has at least $m+1$ non-triangular sectors, while a reduced vertex has at most two.
Therefore $m\ge3$ is impossible: in a reduced arrangement all vertices of multiplicity $\ge3$ are reduced.
\end{proof}

In particular, the pruning of non-reduced arrangements can be combined with the first-letter fixation, and then only $(0,2)$ remains.

With the pruning of non-reduced arrangements, the search uses only letters $(g,m)$ with $m\le4$ and discards a branch as soon as the known non-triangular sectors of some vertex of multiplicity 3 or 4 contradict \cref{def:reduced}.
A sector counts as non-triangular once its face has closed as a non-triangle or has become doomed (\cref{app:doomed}). The other sectors are either actual triangles or can still become triangles.
This rule discards no words of reduced arrangements.
Under this restriction the search no longer emits all words.
Nevertheless, the output contains all words of maximal length, so it is empty if and only if no arrangement with $a_3\ge K$ exists.
If the state table pruning is used as well, its entries must also store the word length (\cref{lem:state-table-reduced}).

\subsection{Computational results}

We have described the search algorithm and justified its correctness. Its implementation, run instructions, output words and run logs are available in the public repository~\cite{nonsimple-arrangements}. It remains to describe the results.

For convenience, we group the restrictions and prunings into four modes (\cref{tab:modes}), each extending the previous one.
Each further mode shrinks the search tree but relies on more proofs, and in modes III and IV the search no longer emits all words.
A word has one letter per vertex, so the longest words represent the arrangements with the largest number of vertices; mode IV still emits all the longest words of mode III.
Running several modes lets us cross-check the outputs and find the shortest chain of proofs behind each result.

\begin{table}[tp]
\centering
\small
\begingroup
\setlength{\tabcolsep}{4pt}
\caption{Search modes. The search emits words with $a_3\ge K$ only; the last column shows the subset guaranteed to be emitted.}
\label{tab:modes}
\renewcommand{\arraystretch}{1.2}
\begin{tabular}{l>{\raggedright\arraybackslash}p{5.3cm}>{\raggedright\arraybackslash}p{2.6cm}>{\raggedright\arraybackslash}p{2.65cm}}
\toprule
mode & restrictions and prunings & justification & guaranteed output \\
\midrule
I & canonical form, budget pruning (P1), early non-triangle count (P2), forced~descent~(P3) &\cref{lem:canonical}, \cref{cor:budget}, Lemmas~\ref{lem:bounded-face},~\ref{lem:forced-descent} & all canonical words \\
II & I + segment pruning (P4), state~table~pruning~(P5), doomed-face count (P6) &\cref{app:prunings}, \ref{lem:segment-balance}--\ref{lem:state-table} & same as I \\
III & II + first-letter fixation ($g_0=0$) & \cref{lem:first-letter} & a word of each arrangement \\
IV & III + pruning of non-reduced arrangements (P7) &Lemmas~\ref{lem:perturbation}, \ref{lem:state-table-reduced}, \cref{cor:reduced} & all the longest words of mode III \\
\bottomrule
\end{tabular}
\endgroup

\bigskip
\caption{Search runs. For $n=11$ ($K=33$) and $n=14$ the target $K$ exceeds the best known $a_3$ by one, and the output is empty; in the other rows $K$ is the maximum.
In mode IV the total number of words depends on the state table size, but the number of the longest words does not.
$^*$Parallel run on 32 vCPUs with a shared state table; its node count is nondeterministic.}
\label{tab:runs}
\bigskip
\begin{tabular}{lllrrrr}
\toprule
$n$ & $K$ & mode & words & longest words & nodes & time \\
\midrule
6 & 7 & I & 112 & 32 & 10\,222 & $<$1 s \\
 & & II & 112 & 32 & 3\,000 & $<$1 s \\
 & & III & 40 & 10 & 1\,028 & $<$1 s \\
 & & IV & 22 & 10 & 643 & $<$1 s \\
\midrule
8 & 15 & I & 32 & 32 & 8\,321\,398 & $<$1 s \\
 & & II & 32 & 32 & 203\,056 & $<$1 s \\
 & & III & 12 & 12 & 74\,465 & $<$1 s \\
 & & IV & 12 & 12 & 48\,619 & $<$1 s \\
\midrule
10 & 25 & I & 2\,848 & 1\,408 & 79\,036\,703\,322 & 41 min \\
 & & II & 2\,848 & 1\,408 & 61\,604\,535 & 12 s \\
 & & III & 1\,002 & 408 & 23\,781\,671 & 5 s \\
 & & IV & 728 & 408 & 16\,326\,051 & 4 s \\
\midrule
11 & 32 & II & 220 & 220 & 597\,776\,192 & 1.9 min \\
 & & III & 80 & 80 & 229\,384\,410 & 42 s \\
 & & IV & 80 & 80 & 157\,657\,333 & 42 s \\
\midrule
11 & 33 & I + ($g_0=0$) & 0 & 0 & 627\,129\,383\,709 & 5.8 h \\
 & & II & 0 & 0 & 204\,761\,604 & 33 s \\
 & & III & 0 & 0 & 78\,929\,213 & 12 s \\
 & & IV & 0 & 0 & 51\,840\,705 & 12 s \\
\midrule
12 & 38 & II & 912 & 384 & 38\,999\,559\,942 & 2.0 h \\
 & & III & 360 & 148 & 11\,972\,835\,436 & 38 min \\
 & & IV & 360 & 148 & 8\,931\,540\,630 & 39 min \\
\midrule
14 & 55 & IV & 0 & 0 & 6\,641\,676\,901\,608\rlap{$^*$} & 22.6 h $\times$ 32\rlap{$^*$} \\
\bottomrule
\end{tabular}
\end{table}

Modes I and II emit identical sets of words for $n=6,8,10$. The gain from the prunings of \cref{app:prunings} grows and reaches a factor of 1283 in the number of tree nodes for $n=10$.
Without mode II the case $n=12$ is out of reach: extrapolating the speed of mode I gives months.
However, as $n$ increases, the growth of the tree should make the state table less efficient, since its entries are constantly evicted.

In modes III and IV, the first-letter fixation $g_0=0$ shrinks the tree by another factor of 2.6--3.3, and the pruning of non-reduced arrangements by a factor of 1.3--1.6.
For $n\le12$ mode IV gains almost nothing in running time, because checking reduced vertices makes each node more expensive, but sampled estimates of subtree sizes for $n=14$ predict a severalfold gain.

The run for $n=14$ used 32 vCPUs of a cloud server.
The task was split by 8-letter word prefixes into 107\,732 subtrees, and their parallel traversal with a shared 32~GB state table took about 22.6 hours (700 vCPU-hours in total).
The other runs in \cref{tab:runs} were performed single-threaded on a laptop with a 4~GB state table.\footnote{The server has an AMD EPYC 9845 processor, and the laptop an AMD Ryzen AI 9 HX 370. In single-threaded runs on a partially loaded server, a server vCPU was slower than a laptop core by a factor of 1.4 in mode I, where the speed is limited by the processor, and by a factor of 1.15 in modes II--IV, where it is limited by memory access.}

The node counts of single-threaded runs are reproducible.
In the parallel run for $n=14$, the total number of nodes (about $6.6\cdot10^{12}$, \cref{tab:runs}) depends on the order in which the parallel traversal fills the shared state table, so it is nondeterministic. This does not affect the verdict (\cref{rem:shared-table}, which carries over to \cref{lem:state-table-reduced}).

\begin{theorem}\label{thm:small-n}
For even $6\le n\le14$, the maximum number of triangles in an arrangement of $n$ lines with vertices of any multiplicity equals the conjectured bound \eqref{eq:conj-bound}, and the same holds for pseudolines:
\begin{equation}\label{eq:small-n}
\begin{array}{c|ccccc}
n & 6 & 8 & 10 & 12 & 14\\ \hline
K(n)=K_{\mathrm{ps}}(n) & 7 & 15 & 25 & 38 & 54
\end{array}
\end{equation}
\end{theorem}

\begin{proof}
\emph{Upper bound.} The search (\cref{tab:runs}) rules out pseudoline arrangements with $a_3\ge K(n)+1$.
For $n=6,8,10$ this is shown by mode I, and for $n=12$ by mode II: these modes emit all canonical words with $a_3\ge K(n)$, and all emitted words have $a_3=K(n)$.
For $n=14$ this is shown by mode IV: with the target $a_3\ge55$ its output is empty, while by \cref{lem:state-table-reduced} it would be nonempty if such an arrangement existed.
Hence $K_{\mathrm{ps}}(n)$ does not exceed the values \eqref{eq:small-n}, and by \eqref{eq:K-ps} neither does $K(n)$.

\emph{Lower bound.} The arrangements of 8 and 14 lines with 15 and 54 triangles are the members of the series of \cref{thm:series} for $t=0$ and $t=1$.
Line arrangements for $n=6,10,12$ are shown in \cref{app:figures} (\cref{fig:n6-all,fig:n10-all,fig:n12-all}). Some of them were known before (for $n=10,12$ see~\cite{khodulev99}).

The cases $n=6,8,10$ use only mode I and depend neither on \cref{app:prunings} nor on \cref{lem:perturbation,lem:first-letter}. The case $n=12$ depends on \cref{app:prunings}, and the case $n=14$ also on \cref{cor:reduced}, which rests on both lemmas.
\end{proof}

The algorithm does not use the parity of $n$ in any way.
For odd $n=5,7,9$ it reproduces the optimal arrangements of our previous exhaustive enumeration~\cite[Section~7]{pu26b}.
We also checked the odd value $n=11$.
It is of interest because the known arrangements of 11 lines have at most 32 triangles, while the bound \eqref{eq:a3-bound} gives $a_3(\mathcal A)\leq33$.
For odd~$n\equiv3,5\pmod 6$, perfect arrangements (those attaining the upper bound \eqref{eq:a3-bound} without rounding) exist up to $n=93$, with $n=11$ the only known exception~\cite[Tables~1 and~4]{pu26b}.
To our knowledge, the absence of arrangements with $a_3=33$ has previously been checked only for simple arrangements~\cite{blanc11,savchuk25}.
Our exhaustive search covers vertices of arbitrary multiplicity.

\begin{theorem}\label{thm:n11}
No arrangement of 11 pseudolines with vertices of any multiplicity has $a_3\ge33$. Consequently, $K(11)=K_{\mathrm{ps}}(11)=32$.
\end{theorem}

\begin{proof}
\emph{Upper bound.} The search in mode I with the first-letter fixation $g_0=0$ and the target $a_3\ge33$ emitted no words (\cref{tab:runs}).
By \cref{lem:first-letter}, every arrangement is represented by a canonical form with $g_0=0$, and mode I emits all such words with $a_3\ge K$.
Hence there are no arrangements of 11 pseudolines with $a_3\ge33$, and by \eqref{eq:K-ps} no such line arrangements either.

\emph{Lower bound.} Arrangements of 11 lines with 32 triangles are known, for example the one found by Kabanovich~\cite{khodulev99}.

The proof relies on \cref{lem:first-letter} but not on \cref{app:prunings} or \cref{lem:perturbation}.
\end{proof}

Modes II--IV give the same empty output about three orders of magnitude faster than mode I (\cref{tab:runs}).
With the target $a_3\ge32$, mode II emits 220 words representing five arrangements, and all of them are simple, as are the optima found for $n=5,7,9$. So for odd $n\le11$ all optima are simple.

\subsection{Optimal arrangements}

Passing from words to arrangements with multiple points requires extending to the non-simple case the $O$-matrix~\cite[Definition~5.1]{pu26b}, which encodes wiring diagrams (cf.~Savchuk's table notation~\cite{savchuk25}).
The bijection between commutation classes of words and $O$-matrices~\cite[Lemma~5.2]{pu26b} and the bijection between $O$-matrices and wiring diagrams representing arrangements~\cite[Remark~5.3]{pu26b} have to be extended as well.
We adapted our deduplication and symmetry analysis programs~\cite[Section~5]{pu26b} to arrangements with multiple points and present their results as observations.
The formal extension requires no new ideas but is lengthy, and we omit it.

The same arrangement is represented by many wiring diagrams, since the direction of the sweep is arbitrary, and each diagram is encoded by a commutation class of words (see~\cite[Observation~2.1]{pu26b}).
A labeling of the pseudolines is determined by three choices: which pseudoline is $0$, an end of it, and a direction around infinity.
An affine arrangement has $4n$ such labelings.
Each of them gives a wiring diagram, and only one word of its commutation class is in canonical form.
An arrangement with symmetry group $H_E$~\cite[Definition~5.11]{pu26b} gives $4n/|H_E|$ canonical words.

To eliminate duplicate words that encode the same arrangement, we repeat the deduplication procedure of~\cite[Subsection~7.2]{pu26b}.
Each word gives an $O$-matrix, to which all Euclidean rotations and reflections and all projective rotations are applied.
The lexicographic minima of the resulting $O$-matrices are canonical representatives of the Euclidean and projective equivalence classes ($E$- and $P$-classes)~\cite[Definitions~5.7 and~5.8]{pu26b}. The result is shown in \cref{tab:symmetries}.

\begin{table}[tp]
\centering
\small
\setlength{\tabcolsep}{4pt}
\caption{Optimal arrangements and their symmetries. $\#E$ and $\#P$ are the numbers of Euclidean and projective classes, $H_E$ is the symmetry group of an affine arrangement, $G_P$ that of its projective closure~\cite[Definition~5.10]{pu26b}; $4C_1$ means four classes with group $C_1$.}
\label{tab:symmetries}
\begin{tabular}{lrrrlrll}
\toprule
$n$ & words & words, $g_0=0$ & $\#E$ & $H_E$ & $\#P$ & $G_P$ & $\sum_E 4n/|H_E|$ \\
\midrule
6 & 112 & 40 & 6 & $4C_1,\ 2C_3$ & 6 & $2C_1,\ 2C_2,\ 2C_3$ & $24\cdot4+8\cdot2$ \\
8 & 32 & 12 & 1 & $C_1$ & 1 & $C_1$ & $32\cdot1$ \\
10 & 2\,848 & 1\,002 & 72 & $71C_1,\ C_5$ & 70 & $63C_1,\ 6C_2,\ C_5$ & $40\cdot71+8\cdot1$ \\
12 & 912 & 360 & 19 & $19C_1$ & 19 & $19C_1$ & $48\cdot19$ \\
\bottomrule
\end{tabular}
\end{table}

Each $E$-class turns out to have $4n/|H_E|$ representatives\footnote{In the odd case the number of representatives was half as large, $2n/|H_E|$~\cite[Lemma~5.13]{pu26b}: for odd $n$ the reflection $g\mapsto n-2-g$ changes the parity of the generators $(g,2)$, and by construction the search of~\cite[Section~4]{pu26b} emitted only words in which odd generators are in the majority.}, that is, an $E$-class is exactly an affine arrangement with unlabeled pseudolines. The sums of the numbers of representatives over all $E$-classes are 112, 32, 2848 and 912, and, as they should, they equal the total numbers of emitted words.
Words with first letter $g_0=0$ form only part of the output, but by \cref{lem:first-letter} they represent all arrangements, and their deduplication gives the same $E$-classes.

Only three optima have affine symmetries: two with $H_E=C_3$ for $n=6$ (\cref{fig:n6-all}) and one with $H_E=C_5$ for $n=10$ (\cref{fig:n10-all}).
The projective closure can be more symmetric: for $n=6$ and $n=10$ some arrangements without affine symmetries have $G_P=C_2$.
For $n=10$, two projective classes each contain two optimal affine arrangements, so the 72 $E$-classes give 70 $P$-classes.

Apart from symmetries, the optimal arrangements differ in their multiple points.

\begin{table}[H]
\centering
\small
\setlength{\tabcolsep}{4pt}
\caption{Multiple points of the optima. All of them are triple points. The table gives the number of $E$-classes with a given number $t_3$ of triple points.}
\label{tab:triple-points}
\begin{tabular}{lcccccccccccc}
\toprule
$t_3$ & 0 & 1 & 2 & 3 & 4 & 5 & 6 & 7 & 8 & 9 & 10 & total \\
\midrule
$n=6$ & 2 & 1 & 1 & 2 &  &  &  &  &  &  &  & 6 \\
$n=8$ &  &  & 1 &  &  &  &  &  &  &  &  & 1 \\
$n=10$ & 36 & 8 & 3 & 5 & 3 & 3 & 1 & 4 & 3 & 2 & 4 & 72 \\
$n=12$ &  &  & 8 & 4 & 5 & 1 & 1 &  &  &  &  & 19 \\
\bottomrule
\end{tabular}
\end{table}

All multiple points of the optima are triple points, and their number varies widely: for $n=10$ it takes all values from 0 to 10 (\cref{fig:n10,fig:n10-all}). Simple optima exist for $n=6$ and $n=10$, where \eqref{eq:blanc-even} coincides with \eqref{eq:conj-bound}. For $n=8$ the bound \eqref{eq:blanc-even} already rules them out, since $8\equiv 2\pmod 6$, and the unique optimum, which has no nontrivial symmetries (\cref{tab:symmetries}), contains exactly two triple points.

For $n=12$ the bound \eqref{eq:blanc-even} also allows 38 triangles, but none of the 19 optima is simple. In the next section we consider ways to construct arrangements of an even number of pseudolines by adding one pseudoline. In particular, we recall that a pseudoline can be added to any perfect arrangement of $n-1$ pseudolines so that the result is a simple arrangement attaining \eqref{eq:blanc-even}. As noted above, perfect arrangements are known for all $n-1\equiv3,5\pmod6$ up to 93, except $n-1=11$. Hence, among $n\equiv0,4\pmod6$, $n=12$ is the only value in the range $n\le94$ without a simple arrangement attaining \eqref{eq:blanc-even}.
